\section{Knowledge Graphs and Link Prediction}\label{sec:fund_kg}

A knowledge graph represents knowledge as a network of entities and their
relationships. Entities correspond to distinguishable real-world objects
or abstract concepts and are represented as nodes, while relations are
represented as typed edges connecting these nodes. Unlike conventional
data representations that treat records independently, a knowledge graph
explicitly captures how different objects relate to one another, providing
a structured and semantically meaningful representation of interconnected
information~\cite{ji2021survey}.

Formally, a knowledge graph can be defined as a collection of facts
represented as triples
\[
  (h, r, t),
\]
where $h$ denotes the head entity, $r$ the relation, and $t$ the tail
entity~\cite{ji2021survey, wang2017kgesurvey}. The relation specifies the
meaning of the connection between the two entities. For example, the
following triples illustrate this structure:
\[
\begin{aligned}
  (\text{Paris},     \;&\text{CAPITAL\_OF},  \text{France})   \\
  (\text{Component},\;&\text{SUPPLIED\_BY},\text{Supplier})  \\
  (\text{Incident},  \;&\text{HAS\_DEFECT},  \text{Defect})
\end{aligned}
\]
The first states that Paris is the capital of France; the second, that a
component originates from a particular supplier; and the third records
the association between a reported incident and its observed defect.

A knowledge graph is formed by combining many such triples, allowing the
same entity to participate in multiple relations. A single component, for
instance, may simultaneously be connected to a supplier, a production
plant, a recorded incident, and an observed defect. The resulting
structure is therefore a multi-relational graph, in which different edge
types express different semantic relationships~\cite{wang2017kgesurvey}.
Entities and relations may additionally carry types and properties: an
incident entity may be described by attributes such as an identifier, a
creation date, and a status, while a defect entity may be described by its
type, location, and affected area. These descriptive properties
characterize individual entities, whereas relations express how entities
are connected; together, they allow heterogeneous information to be
represented within a common structure while preserving its original
domain-specific meaning~\cite{ji2021survey}.

This representation is particularly suitable for domains in which
relevant information is distributed across multiple sources. In
industrial quality management, for example, analyzing a reported incident
may require information about defects, materials, plants, suppliers,
inspections, rework, and containment actions. These objects describe
different aspects of the same underlying context; representing them as
connected entities makes it possible to retrieve relevant information
through their relationships, rather than examining each source
independently. The same principle applies more broadly: wherever an
organization's knowledge is distributed across heterogeneous records that
share common objects, a knowledge graph offers a structured way to
reconnect that information for retrieval and reasoning.

\subsection{Knowledge Graph Completion and Link Prediction}\label{sec:fund_kg_completion}

Knowledge graphs are generally incomplete, because not every valid
relationship is observed or explicitly recorded. The absence of a link
does not necessarily imply that no relationship exists; it may instead
indicate that the corresponding fact has not yet been identified, entered,
or derived from the available information. Knowledge graph completion
addresses this incompleteness by identifying plausible missing facts from
the existing graph structure~\cite{ji2021survey}.

Link prediction is one of the principal tasks used for this purpose.
Given an incomplete triple, link prediction estimates the missing entity
or evaluates whether a proposed relationship is
plausible~\cite{wang2017kgesurvey, rossi2021comparative}. The task may
take different forms:
\[
\begin{aligned}
  (h, r, ?) &\quad \text{the tail entity is unknown}, \\
  (?, r, t) &\quad \text{the head entity is unknown}, \\
  (h, r, t) &\quad \text{the plausibility of a complete triple is evaluated}.
\end{aligned}
\]
For example, given the incomplete triple
\[
  (\text{Incident}, \text{CAUSED\_BY}, ?),
\]
the link prediction task is to identify the most plausible causer among a
set of candidate entities. Candidates are typically ranked according to a
scoring function that estimates the plausibility of each possible
completion; the highest-ranked candidate is the entity considered most
compatible with the learned graph patterns. Link prediction therefore
transforms the graph completion problem into a structured prediction task
over entities and relations~\cite{rossi2021comparative,
wang2017kgesurvey}.

Knowledge graph embedding methods commonly support this task by mapping
entities and relations into a continuous vector space. A scoring function
then evaluates how plausible a given combination of head entity, relation,
and tail entity is; the representations are learned such that observed
triples receive higher plausibility scores than incorrect or unobserved
combinations. However, the ability of such a model to handle new entities
depends on whether it follows a transductive or an inductive learning
setting.

\subsection{Transductive and Inductive Link Prediction}\label{sec:fund_kg_transind}

In a \textbf{transductive} setting, the set of entities considered during
inference is already known at training time: the model observes a graph
containing all relevant entities and learns representations for them,
even if some of their relationships are hidden and later used for
evaluation. Conventional knowledge graph embedding methods frequently
follow this paradigm by assigning a trainable embedding vector to each
entity in the training graph~\cite{wang2017kgesurvey, teru2020grail}. Such
models can predict missing links between entities for which embeddings
have already been learned; for example, if an incident and all of its
possible causers were included in the training graph, the model could
estimate which of the missing relations between these known entities is
most plausible. This setting, however, assumes that the entities required
during inference were already available at training time: an
entity-specific embedding cannot be retrieved for an entity that did not
occur in the training graph~\cite{teru2020grail, rossi2021comparative}.

This limitation is significant in dynamic application environments, where
new products, incidents, materials, or other business objects may
continuously be added after a model has been trained. Retraining the
entire model whenever a new entity appears is generally undesirable, and,
more importantly, a model that depends only on memorized entity
representations does not demonstrate that it has learned relational
patterns transferable to previously unseen cases.

\textbf{Inductive} link prediction addresses these limitations by
requiring a model to generalize to entities that were not part of the
training graph~\cite{teru2020grail, galkin2022ilpc}. During inference, the
graph may contain entities absent from training; instead of relying
exclusively on a stored embedding for every entity, an inductive method
derives information from transferable features, relations, or structural
patterns in an entity's graph neighborhood. A neighborhood describes the
entities and relations surrounding a target entity; for an incident, this
neighborhood may include its associated defect, material, plant, and
inspection or rework records. Although the incident itself is new, some
of these neighboring entity types and relational structures may resemble
patterns observed in historical incidents, and an inductive method can
use the structure surrounding the new entity to support its prediction.
Subgraph-based inductive approaches formalize this idea by extracting a
local subgraph around the candidate entities and learning the relational
and structural patterns within such subgraphs, rather than only the
identities of individual entities~\cite{teru2020grail}; the learned
patterns can then be applied to a new subgraph containing an entity that
was absent during training.

The distinction can be summarized as follows: transductive link
prediction predicts missing relations among entities that are already
known during training, whereas inductive link prediction transfers
learned relational patterns to entities that were not present during
training. This distinction concerns the availability of \emph{entities},
rather than only the presence or absence of individual edges: hiding a
relation for an incident that was already observed during training
remains a transductive task, whereas predicting the same type of relation
for an incident that was entirely absent from training constitutes an
inductive task.

\subsection{Relevance to the Proposed Methodology}\label{sec:fund_kg_relevance}

RQ1 (Section~\ref{sec:intro_rq}) asks how heterogeneous industrial
knowledge can be retrieved through a knowledge graph to support the
investigation of \emph{previously unseen} quality incidents. The concepts
introduced above give this question a precise formulation: causer
attribution is framed as the prediction of a missing relation between a
quality incident and a candidate causer,
\[
  (q, \text{CAUSED\_BY}, c),
\]
where $q$ denotes a quality incident and $c$ a candidate causer.
Historically resolved incidents provide examples for which this relation
is already known; a newly created incident has no known causer and must
instead serve as the head entity of an incomplete triple,
\[
  (q_{\text{new}}, \text{CAUSED\_BY}, ?).
\]

Because $q_{\text{new}}$ did not exist when the historical training graph
was built, no entity-specific representation could have been learned for
it -- this is precisely the inductive setting described in
Section~\ref{sec:fund_kg_transind}. What remains available is
$q_{\text{new}}$'s local graph context: its connections to material,
plant, inspection, and other quality information. Rather than depending
on the identity of the incident, the methodology developed in this thesis
compares this relational neighborhood against the structures of
previously resolved cases, directly operationalizing RQ1's requirement to
retrieve evidence for incidents that have no exact historical match.

This inductive framing is also what connects the proposed methodology to
the broader scientific motivation outlined in
Section~\ref{sec:intro_motivation_sci}: treating causer attribution as
link prediction under an open, continuously growing entity set, rather
than a fixed, previously observed one, is itself the methodological
contribution that generalizes beyond BMW's quality process to any
industrial decision-support setting in which an organization's knowledge
graph keeps growing after deployment. Chapter~\ref{ch:setup} reflects
this requirement directly: evaluation holds out incidents whose entities
are entirely unseen during training, since a transductive evaluation
would not reflect this intended, continuously evolving setting.
